% J2.1: flujo de un proyecto de análisis; lo crudo no se toca, lo regenerable se puede borrar.
\documentclass[tikz,border=4pt]{standalone}
\usepackage{resonancia}
\begin{document}
\begin{tikzpicture}
  \node[cajaAzul, minimum width=3.0cm] (cr) at (0,0) {\textbf{datos/crudos/}\\[-1pt]{\footnotesize no se edita nunca}};
  \node[cajaVerde, minimum width=3.0cm] (sc) at (4.6,0) {\textbf{scripts/}\\[-1pt]{\footnotesize \texttt{run\_all.jl}}};
  \node[cajaNaranja, minimum width=3.4cm] (re) at (9.4,0.9) {\textbf{resultados/tablas, figuras}\\[-1pt]{\footnotesize regenerables}};
  \node[caja, draw=rVioleta, fill=rVioleta!8, minimum width=3.4cm] (lo) at (9.4,-0.9) {\textbf{resultados/logs}\\[-1pt]{\footnotesize qué se hizo y cuándo}};
  \node[caja, draw=rGris, fill=rGrisClaro!50, minimum width=3.0cm] (no) at (4.6,-1.9) {\textbf{notas/}\\[-1pt]{\footnotesize decisiones y dudas}};
  \draw[flecha] (cr) -- (sc); \draw[flecha] (sc) -- (re); \draw[flecha] (sc) -- (lo);
  \draw[flecha, color=rGris, dashed] (sc) -- (no);
  \node[etiqueta, color=rGris, anchor=north] at (4.6,-2.6) {el entorno (\texttt{Project.toml}, \texttt{Manifest.toml}) y la semilla completan la receta};
\end{tikzpicture}
\end{document}
